
% !TEX root = ../../../Esperimenti/.tex/MEF.tex
% LTeX: language=it

% La sfumatura degli archi approssimati/rilassati usa «path fading», che
% richiede la libreria «fadings».
\usetikzlibrary{fadings}

\begin{tikzpicture}%
  \nodeDiagramSettings

  %\begingroup Coordinate dei vertici
    \newcommand\vertexCoordinates{
      \coordinate (v1) at (0,0);
      \coordinate (v2) at (2,2);
      \coordinate (v3) at (-1,2.5);
      \coordinate (v4) at (1.5,4.0);
      \coordinate (v5) at (3.0,-0.5);
      \coordinate (v6) at (4,3);
      \coordinate (v7) at (5.0,0.5);
    }

    % \newcommand\vertexCoordinates{
    %   \coordinate (v1) at (0,0);
    %   \coordinate (v2) at (2.5,1.5);
    %   \coordinate (v3) at (-1,2.5);
    %   \coordinate (v4) at (1.5,4.0);
    %   \coordinate (v5) at (3.0,-0.5);
    %   \coordinate (v6) at (4.5,2.8);
    %   \coordinate (v7) at (5.0,0.5);
    % }

    \newcommand\drawEmptyNodes[1]{
      \vertexCoordinates
      \foreach \i in {1,2,...,7}
        \node[
          nodeStyle,
          numberStyle,
          fill=white
        ] (#1\i) at (v\i) {\i};

      % \begin{scope}[on background layer]
        \foreach \i in {1,2}
          \node[starStyle,rotate=\i*45] at (v\i) {};
      % \end{scope}
    }

    \newcommand\drawUniformedNodes[2]{
      \vertexCoordinates
      \foreach \i in {1,2,...,7}
        \node[nodeStyle,fill=#2] (#1\i) at (v\i) {};

      % \begin{scope}[on background layer]
        \foreach \i in {1,2}
          \node[
            starStyle,
            % fill=cStar,
            % fill opacity=0.9,
            rotate=\i*45
          ] at (v\i) {};
      % \end{scope}
    }

    \newcommand\drawColoredNodes[1]{
      \vertexCoordinates
      \foreach \i in {1,2,...,7}
        \node[nodeStyle,fill=color\i] (#1\i) at (v\i) {};

      % \begin{scope}[on background layer]
        \foreach \i in {1,2}
          \node[
            starStyle,
            fill=cStar,
            % fill opacity=0.9,
            rotate=\i*45
          ] at (v\i) {};
      % \end{scope}
    }

    % \newcommand\drawColoredNodes[1]{\vertexCoordinates
    %   \node[nodeStyle,fill=cTeal] (#11) at (v1) {};
    %   \node[nodeStyle,fill=cPurple] (#12) at (v2) {};
    %   \node[nodeStyle,fill=cRed] (#13) at (v3) {};
    %   \node[nodeStyle,fill=cOrange] (#14) at (v4) {};
    %   \node[nodeStyle,fill=cRed] (#15) at (v5) {};
    %   \node[nodeStyle,fill=cTeal] (#16) at (v6) {};
    %   \node[nodeStyle,fill=cPurple] (#17) at (v7) {};
    % }
  %\endgroup

  %\begingroup Lati del grafo completo
    % Versione precedente (grafo completo con tutti i 21 spigzi), conservata
    % per riferimento.
    % \newcommand\drawCompleteEdges[1]{
    %   \begin{scope}[on background layer]
    %     \foreach \i [evaluate=\j as \s using int(\i+1)] in {1,...,6}
    %       \foreach \j in {\s,...,7}{
    %         \ifcase\i % Caso 0
    %         \or % Caso 1
    %           \ifnum\j=2\def\b{15}
    %           \else\ifnum\j=5\def\b{-20}
    %           \else\ifnum\j=6\def\b{-10}
    %           \else\ifnum\j=7\def\b{-10}
    %           \else\def\b{20}
    %           \fi\fi\fi\fi
    %         \or % Caso 2
    %           \ifnum\j=3\def\b{-10}
    %           \else\ifnum\j=4\def\b{0}
    %           \else\ifnum\j=5\def\b{-15}
    %           \else\ifnum\j=6\def\b{0}
    %           \else\ifnum\j=7\def\b{10}
    %           \else\def\b{20}
    %           \fi\fi\fi\fi\fi
    %         \or % Caso 3
    %           \ifnum\j=5\def\b{-20}
    %           \else\ifnum\j=7\def\b{-10}
    %           \else\def\b{20}
    %           \fi\fi
    %         \or % Caso 4
    %           \def\b{20}
    %         \or % Caso 5
    %           \def\b{-20}
    %         \or % Caso 6
    %           \def\b{20}
    %         \else % Caso 7
    %           \def\b{20}
    %         \fi
    %         \ifnum\i=1
    %           \draw[nodeArrowApprox] (#1\i) to[bend left=\b] (#1\j);
    %         \else
    %           \draw[nodeArrowApprox,transparentNodeArrow] (#1\i) to[bend left=\b] (#1\j);
    %         \fi
    %       }
    %   \end{scope}
    % }
  %\endgroup

  %\begingroup Lati con topologia incognita (grafo approssimato/rilassato)
    % La topologia non e' nota: ogni vertice conserva il grado che ha nel
    % grafo esatto, ma gli spigzi non raggiungono alcun nodo in particolare.
    % Sono archi sfumati, curvati come gli spigzi originali e orientati verso
    % l'interno dell'inviluppo convesso {1,3,4,5,6,7}, in direzioni casuali:
    % sfumano invece di essere troncati a meta'.
    %
    % Gradi nel grafo esatto (spigzi 1-2, 2-4, 2-5, 2-6, 3-4, 5-7, 6-7):
    %   deg(1)=1, deg(2)=4, deg(3)=1, deg(4)=2, deg(5)=2, deg(6)=2, deg(7)=2
    % (da mantenere allineati con la lista nel grafo esatto).
    %
    % Ogni terna e': vertice / angolo in gradi / curvatura / lunghezza
    % dell'arco (curvatura positiva = curva a sinistra, come «bend left»).
    % Le direzioni sono state scelte in modo che ogni arco termini
    % nell'inviluppo convesso, non passi sopra alcun nodo (distanza minima
    % 0.86 dal centro, contro un raggio di 0.35) e non si sovrapponga agli
    % altri archi (distanza minima 0.22). Eventuali sovrapposizioni vanno
    % risolte cambiando direzione (o accorciando l'arco).
    %
    % La sfumatura e' data da «path fading=east» (richiede la libreria
    % «fadings»), orientato con «fading angle»: opaco sul vertice, trasparente
    % sulla punta.
    \newcommand\drawApproxEdges[1]{
      \begin{scope}[on background layer]
        \foreach \i/\ang/\bnd/\dst in {
          1/35/25/1.2,
          3/-33.5/25/1.5,
          4/-113/30/1.05,
          4/-59.6/30/1.35,
          5/110.1/-25/1.5,
          5/147.8/30/1.35,
          6/-93.9/-30/1.35,
          6/-150.3/-25/1.2,
          7/125.6/-30/1.05,
          7/-174.3/25/1.35,
          2/61.1/-25/1.2,
          2/151.1/25/1.05,
          2/241.1/-25/1.35,
          2/331.1/25/1.05
        }{
          \draw[nodeArrowApprox, solid, path fading=east, fading angle=\ang]
            (#1\i) to[bend left=\bnd] ++(\ang:\dst);
        }
      \end{scope}
    }
  %\endgroup

  \begin{scope}[
    local bounding box=tikZFigure,
    scale=\plotScale,
    transform shape,
  ]
    %\begingroup Grafo esatto
      % Il grafo e' reso come non orientato: gli spigzi sono linee semplici
      % senza punte perche' la topologia e' indipendente dalle leggi di
      % interazione, che sono un concetto a se'.
      \tikzbox{{(0,0)}}{ExactGraph}{
        % \drawColoredNodes{e}
        % \drawUniformedNodes{e}{none}
        \drawEmptyNodes{e}

        \begin{scope}[on background layer]
          \foreach \i/\j/\b/\t in {
            1/2/20/0,
            2/4/20/1,
            2/5/20/1,
            2/6/20/1,
            3/4/20/1,
            % 3/6/20/1,
            5/7/20/1,
            6/7/20/1
          }{
            \ifnum\t=0
              \draw[indirectArrow] (e\i) to[bend left=\b] (e\j);
            \else
              \draw[indirectArrow,transparentNodeArrow] (e\i) to[bend left=\b] (e\j);
            \fi
          }
        \end{scope}
      }
    %\endgroup

    %\begingroup Grafo approssimato
      \tikzbox[north]{{
        ($(ExactGraph.south east)
        +(.5*\xPlotSep em,-\yPlotSep em)$)
      }}{ApproxGraph}{
        % \drawColoredNodes{a}
        % \drawUniformedNodes{a}{none}
        \drawEmptyNodes{a}
        \drawApproxEdges{a}
      }
    %\endgroup

    %\begingroup Grafo rilassato
      \tikzbox{{
        ($(ExactGraph.south east)
        +(\xPlotSep em,0)$)
      }}{RelaxedGraph}{
        % \drawColoredNodes{r}
        \drawUniformedNodes{r}{white}
        % \drawUniformedNodes{r}{none}
        % \drawEmptyNodes{r}
        \drawApproxEdges{r}
      }
    %\endgroup

    %\begingroup Testo
      \node[titleStyle] at (ExactGraph.north) {\ifcasscclass Exact topology\else Topologia esatta\fi};
      \node[titleStyle] at (ApproxGraph.north) {\ifcasscclass Approximated topology\else Topologia approssimata\fi};
      \node[titleStyle] at (RelaxedGraph.north) {\ifcasscclass Relaxed topology\else Topologia rilassata\fi};
      
      \node[commentStyle,below=0.5em of e1] {$(i,s)$};
      \node[commentStyle,below=0.5em of a1] {$(i,s)$};
      \node[commentStyle,below=0.5em of r1] {$(x,s)$};
    %\endgroup

    %\begingroup Frecce
      \draw[transformingArrow] 
        ($(ExactGraph.south west)!0.2!(ExactGraph.south east)$) 
        to[bend right=40] node[arrowCommentStyle]
        {\ifcasscclass\relax\else\cref{defMatriceGradi}\fi}
        % {Def. \ref{defMatriceGradi}}
        % {Def. \zref[default]{defMatriceGradi}}
        % {Def. 3.12}
        (ApproxGraph.west);

      \draw[transformingArrow]
        (ApproxGraph.east) to[bend right=40] node[arrowCommentStyle] 
        {\ifcasscclass\relax\else\cref{teoRilassamentoTopologia}\fi}
        % {Teo. \ref{teoRilassamentoTopologia}}
        % {Teo. \zref[default]{teoRilassamentoTopologia}}
        % {Teo. 3.1}
        ($(RelaxedGraph.south west)!0.8!(RelaxedGraph.south east)$);

      % \draw[transformingArrow] (ExactGraph.south) to[bend right=40] (ApproxGraph.west);
      % \draw[transformingArrow] (ApproxGraph.east) to[bend right=40] (RelaxedGraph.south);

      % \draw[transformingArrow] (ExactGraph.south) |- (ApproxGraph.west);
      % \draw[transformingArrow] (ApproxGraph.east) -| (RelaxedGraph.south);
    %\endgroup

    %\begingroup Scatole di delimitazione
      % \draw[red,dashed] (ExactGraph.south west) rectangle (ExactGraph.north east);
      % \draw[blue,dashed] (ApproxGraph.south west) rectangle (ApproxGraph.north east);
      % \draw[green,dashed] (RelaxedGraph.south west) rectangle (RelaxedGraph.north east);
    %\endgroup
  \end{scope}

  %\begingroup Riferimento
    % \begin{pgfonlayer}{foreground}
    %   \node[commentStyle] (origine) at (0,0) {origine};
  
    %   \node[commentStyle] (NW) at (ExactGraph.north west) {NW};
    %   \node[commentStyle] (NE) at (ExactGraph.north east) {NE};
    %   \node[commentStyle] (SW) at (ExactGraph.south west) {SW};
    %   \node[commentStyle] (SE) at (ExactGraph.south east) {SE};
    % \end{pgfonlayer}
  
    % \begin{scope}[on background layer]
    %   \draw[style=help lines] (tikZFigure.north west) grid[step=1em] (tikZFigure.south east);
    % \end{scope}
  %\endgroup
\end{tikzpicture}
